\documentclass[11pt,a4paper]{article}

% ── Font (pdflatex-safe, modern sans) ──
\usepackage[T1]{fontenc}
\usepackage[utf8]{inputenc}
\usepackage[scaled=0.92]{helvet}  % Helvetica (modern sans)
\renewcommand{\familydefault}{\sfdefault}

% ── Layout ──
\usepackage[
top=28mm, bottom=30mm, left=26mm, right=26mm,
headheight=14pt
]{geometry}
\usepackage{parskip}
\setlength{\parskip}{0.45em}

% ── Packages ──
\usepackage{xcolor}
\usepackage{tikz}
\usepackage{fancyhdr}
\usepackage{titlesec}
\usepackage{microtype}
\usepackage{hyperref}
\usepackage{amssymb,amsmath}

% ── Colours ──
\definecolor{deep}{HTML}{1B2A4A}
\definecolor{accent}{HTML}{C0392B}
\definecolor{steel}{HTML}{5D6D7E}
\definecolor{light}{HTML}{D5DBDB}

% ── Hyperref ──
\hypersetup{
	colorlinks=true,
	linkcolor=accent,
	urlcolor=steel,
	pdfauthor={Carlo Perassi},
	pdftitle={Ten Things You Cannot Afford to Ignore},
}

\title{Ten Things You Cannot Afford to Ignore}
\author{Carlo Perassi}
\date{March 2026}

% ── Section styling ──
\titleformat{\section}
{\Large\bfseries\color{deep}}
{\thesection.}{0.5em}{}
[\vspace{-0.3em}{\color{accent}\rule{\textwidth}{0.6pt}}]

% ── Header / Footer ──
\pagestyle{fancy}
\fancyhf{}
\renewcommand{\headrulewidth}{0pt}
\fancyfoot[C]{\small\color{steel}\thepage}

% ── Item card environment ──
\newenvironment{card}[3]{%
	% #1 = number, #2 = title, #3 = area tag
	\vspace{0.8em}
	\noindent
	\begin{tikzpicture}[baseline=(num.base)]
		\node[fill=deep, text=white, font=\Large\bfseries,
		minimum width=9mm, minimum height=9mm,
		inner sep=0pt, rounded corners=2pt] (num) {#1};
	\end{tikzpicture}%
	\hspace{0.4em}%
	{\Large\bfseries\color{deep}#2}%
	\hfill
	{\small\color{accent}\MakeUppercase{#3}}%
	\par\vspace{0.15em}\noindent\ignorespaces
}{%
	\par\vspace{0.4em}
	\noindent{\color{light}\rule{\textwidth}{0.4pt}}
}

\begin{document}
	
	% ── TITLE PAGE ──
	\thispagestyle{empty}
	\vspace*{2.5cm}
	
	\begin{center}
		{\color{accent}\rule{0.4\textwidth}{1pt}}\\[1.2em]
		{\fontsize{28}{34}\selectfont\bfseries\color{deep}
			Ten Things You Cannot\\[0.15em] Afford to Ignore}\\[1.2em]
		{\color{accent}\rule{0.4\textwidth}{1pt}}\\[2em]
		{\large\color{steel}\itshape
			An inventory of ideas without which\\[0.3em]
			no contemporary thought can call itself contemporary.}\\[3em]
		{\color{deep}\bfseries Carlo Perassi}\\[0.3em]
		{\small\color{steel} Octopus Lab\enspace\textperiodcentered\enspace Kiwifarm Srl}\\[0.3em]
		{\small\color{steel} March 2026}
	\end{center}
	
	\vfill
	
	\noindent\textcolor{steel}{\small
		Despite the widespread availability of communication media, many people
		essentially operate with a culture that is not of the twentieth century
		but of the nineteenth. The following ten landmarks, drawn from
		mathematics, physics, computer science, logic, and the philosophy of
		science, are non-negotiable prerequisites for anyone who aspires to
		formulate original, \emph{contemporary} thought.}
	
	\newpage

\begin{abstract}
\noindent This article collects ten mathematical, scientific, computational,
and philosophical results that form part of the minimum background for
contemporary thought. The emphasis is not encyclopedic coverage, but structural
orientation: each item marks a limit, framework, or conceptual shift without
which modern reasoning becomes distorted.
\end{abstract}

\section*{Introduction}

% ── THE MERGE ALL TOGETHER APPROACH ──
\subsection*{The ``merge all together'' approach}

When synthesised, these ten pillars do not merely coexist; they interlock to form a single, unified epistemology. 

Reality is not a clockwork stage but a dynamic, observer-dependent geometry (\emph{Relativity}) whose foundational layer is probabilistic and entangled (\emph{Quantum Mechanics}). The ultimate currency of this reality is \emph{information} (\emph{Shannon}). 

When we attempt to decode this information, we are immediately met with hard limits. Our formal descriptions cannot prove every truth (\emph{G\"odel}), our algorithms cannot solve every problem (\emph{Turing}), and even systems with perfect deterministic rules defeat long-term prediction (\emph{Chaos}). Furthermore, finding solutions to the problems we \emph{can} compute is often blocked by insurmountable resource walls (\emph{Computational Complexity}).

How, then, do we proceed? By recognising that human knowledge is the search for the most compressed algorithmic description of reality (\emph{Kolmogorov}). Because absolute certainty is structurally denied to us, we must navigate the world using graduated, probabilistic updates based on incoming evidence (\emph{Bayesian Inference}). And we must remain acutely aware that our reigning models are never final truths, but temporary, compressible frameworks destined to be overthrown when anomalies inevitably pile up (\emph{Kuhn}).

To ``merge it all together'' is to accept a profound and humbling worldview: \textbf{we inhabit a universe of irreducible uncertainty and bounded computation, where knowledge is not the accumulation of absolute facts, but the continuous, probabilistic compression of information.}

\vspace{2.5em}

\begin{center}
	\color{accent}\rule{0.25\textwidth}{0.6pt}
\end{center}

\newpage

	% ── ON THE ORDERING ──
	\subsection*{On the ordering}
	
	The sequence is neither chronological nor alphabetical; it follows a
	\emph{logical dependency chain}, each result presupposes, or sharpens,
	the ones before it.
	
	\textbf{Layer 1: the absolute limits of formal reason} (items 1--2).
	G\"odel and Turing establish, respectively, that no single formal system
	can capture all mathematical truth, and that no algorithm can decide
	all well-posed questions. These are meta-results: they constrain every
	subsequent framework. They come first because they delimit the
	\emph{entire playing field}.
	
	\textbf{Layer 2: the physical world is not what you assumed} (items
	3--4). Relativity and quantum mechanics dismantle the Newtonian stage on
	which pre-twentieth-century thought operates. They appear immediately
	after the logical limits because they show that reality
	itself, space, time, matter, causality, is stranger than any
	nineteenth-century formalism could express.
	
	\textbf{Layer 3: information as a first-class concept} (item 5).
	Shannon's information theory bridges physics and computation. It
	depends on the probabilistic worldview opened by quantum mechanics and
	provides the quantitative language without which items 7--10 cannot be
	stated precisely.
	
	\textbf{Layer 4: how knowledge itself changes} (item 6). Kuhn's
	philosophy of science is placed after the physical revolutions it
	\emph{explains}. One needs relativity and quantum mechanics as concrete
	examples before the concept of ``paradigm shift'' acquires its full
	force.
	
	\textbf{Layer 5: reasoning under uncertainty} (item 7). Bayesian
	inference is the operational counterpart of Kuhn's insight: if
	paradigms shift, certainty is unavailable, and rational belief must be
	graduated. This requires Shannon's notion of information to be
	quantitatively precise.
	
	\textbf{Layer 6: the limits of prediction, computation, and
		description} (items 8--10). Chaos theory shows that determinism does
	not imply predictability. Complexity theory shows that computability
	does not imply tractability. Kolmogorov complexity shows that
	describability itself has an irreducible, uncomputable core. These
	three results close the architecture by revisiting, at a finer grain,
	the limits opened by G\"odel and Turing, completing the arc.
	
	\vspace{0.5em}
	\noindent{\color{accent}\rule{\textwidth}{0.6pt}}
	
	\newpage
	
	% ── THE TEN ITEMS ──
	
	\begin{card}{1}{G\"odel's Incompleteness Theorems}{Logic / Mathematics}
		In 1931 Kurt G\"odel proved that any consistent formal system powerful
		enough to express elementary arithmetic contains true statements that
		\emph{cannot} be proved within the system, and that such a system
		cannot prove its own consistency. These results shattered the
		Hilbert programme and placed permanent, structural limits on what
		formal reasoning can achieve. Anyone who speaks of ``proof,''
		``certainty,'' or ``foundations'' without awareness of incompleteness
		is reasoning inside a nineteenth-century frame that was \emph{already
			refuted} a century ago.
	\end{card}
	
	\begin{card}{2}{Turing Computability}{Computer Science / Mathematics}
		Alan Turing's 1936 paper introduced the abstract machine that bears his
		name and proved the existence of well-posed problems that \emph{no
			algorithm can solve}, most famously, whether an arbitrary programme
		will halt. Computability theory defines the outer boundary of what
		machines (and formalised procedures in general) can ever accomplish.
		It is the operational twin of G\"odel's result: not all truths are
		reachable by systematic computation. Without this, one cannot reason
		honestly about algorithms, artificial intelligence, or automation.
	\end{card}
	
	\begin{card}{3}{Special and General Relativity}{Physics / Geometry}
		Einstein's theories (1905, 1915) replaced Newtonian absolute space and
		time with a unified spacetime whose geometry is shaped by mass-energy.
		Time is not universal; simultaneity is observer-dependent; gravity is
		curvature. These are not exotic curiosities, they underpin GPS,
		cosmology, and every modern understanding of causality. A worldview
		that tacitly assumes a Newtonian stage is not merely imprecise: it is
		\emph{wrong} about the very arena in which events occur.
	\end{card}
	
	\begin{card}{4}{Quantum Mechanics}{Physics / Epistemology}
		The quantum revolution (Planck, Bohr, Heisenberg, Schr\"odinger,
		Dirac, 1900--1930) revealed that nature at its most fundamental level
		is probabilistic, that measurement participates in determining
		outcomes, and that the classical notion of a particle with definite
		trajectory is an approximation. Entanglement, superposition, and the
		uncertainty principle are not philosophical thought-experiments: they
		are the basis of every semiconductor, laser, and, imminently, quantum
		computer. Ignoring them means misunderstanding the material substrate of
		modern technology.
	\end{card}
	
\newpage

	\begin{card}{5}{Shannon's Information Theory}{Mathematics / Engineering}
		Claude Shannon's 1948 paper gave the first rigorous, quantitative
		definition of \emph{information}, independent of meaning, and
		established fundamental theorems on channel capacity, compression, and
		error correction. Information theory is the invisible grammar of the
		digital age: it governs how data is stored, transmitted, encrypted, and
		compressed. Without Shannon, one cannot understand why a signal is
		distinct from noise, or why perfect communication over an imperfect
		channel is possible at all.
	\end{card}
	
	\begin{card}{6}{Kuhn's \emph{Structure of Scientific Revolutions}}{Philosophy of Science}
		Thomas Kuhn (1962) argued that science does not progress by linear
		accumulation but through paradigm shifts: long periods of ``normal
		science'' punctuated by crises and revolutionary replacements of the
		reigning framework. The concepts of \emph{paradigm},
		\emph{incommensurability}, and \emph{anomaly} transformed how we
		understand scientific change and, by extension, how we evaluate
		claims of knowledge. Anyone who treats ``science'' as a monolithic,
		ever-advancing edifice has not reckoned with Kuhn.
	\end{card}
	
	\begin{card}{7}{Bayesian reasoning and statistical inference}{Statistics / Epistemology}
		Bayes' theorem (18th century in origin, 20th century in maturation)
		provides a coherent framework for updating beliefs in the light of
		evidence. Modern science, medicine, machine learning, and decision
		theory are \emph{structurally} Bayesian: they deal in degrees of
		belief, prior distributions, and likelihood ratios rather than
		black-and-white ``proof.'' A person who reasons only in terms of
		``true or false'' is using a logic inadequate for a world of partial
		information and uncertain evidence.
	\end{card}
	
	\begin{card}{8}{Chaos and sensitivity to initial conditions}{Mathematics / Physics}
		Poincar\'e glimpsed it; Lorenz made it vivid in 1963. Deterministic
		systems can be \emph{practically} unpredictable because tiny
		differences in initial conditions grow exponentially over time.
		This is not randomness: the equations are exact but it means that
		long-range prediction is structurally impossible for a wide class of
		systems (weather, turbulence, many biological processes). The
		nineteenth-century dream of Laplacian determinism ``give me the
		initial conditions and I will predict the future'' fails not for
		lack of data, but for mathematical reasons.
	\end{card}

\newpage

	\begin{card}{9}{Computational Complexity (P vs NP)}{Computer Science / Mathematics}
		Even among computable problems, not all are \emph{tractable}. The
		theory of complexity classes (P, NP, PSPACE, and beyond) classifies
		problems by the resources required to solve them as input size grows.
		The still-open question of whether $\mathrm{P} = \mathrm{NP}$ is among
		the deepest in mathematics, and its practical consequences are
		immediate: if $\mathrm{P} \neq \mathrm{NP}$ (as widely believed), an
		enormous class of optimisation, scheduling, and verification problems
		have \emph{no} efficient general solution. Ignoring this means
		misunderstanding what software can and cannot do, and why cryptography
		works.
	\end{card}
	
	\begin{card}{10}{Kolmogorov Complexity}{Mathematics / Computer Science}
		Kolmogorov complexity measures the intrinsic information content of an
		object as the length of the shortest programme that produces it.
		It unifies concepts of randomness, compressibility, and structure in a
		single, mathematically rigorous framework, and it is
		\emph{uncomputable}, which links it directly back to G\"odel and
		Turing. A string is ``random'' if and only if it cannot be compressed;
		a scientific law is ``good'' if it compresses the data it explains.
		Without this lens, discussions of simplicity, Occam's razor, and
		pattern recognition remain purely intuitive and pre-formal.
	\end{card}
	
	\vspace{1.5em}
	
	\begin{center}
		\color{accent}\rule{0.25\textwidth}{0.6pt}
	\end{center}
	
	\vspace{0.5em}
	
	\section*{Conclusion}

	\noindent\textcolor{steel}{\small\itshape
		These ten results are not optional cultural enrichment. They are the
		load-bearing walls of contemporary thought. A mind that has not
		absorbed them can still be erudite in a nineteenth-century sense, but
		it cannot produce ideas that are genuinely \emph{of this century}.}

\end{document}
